\documentclass[11pt,reqno]{amsart}

% ============================================================
% Packages
% ============================================================
\usepackage{amsmath,amssymb,amsthm}
\usepackage{mathtools}
\usepackage{enumitem}
\usepackage{booktabs}
\usepackage{array}
\usepackage[margin=1in]{geometry}
\usepackage{hyperref}
\usepackage{xcolor}

% ============================================================
% Theorem environments
% ============================================================
\theoremstyle{plain}
\newtheorem{theorem}{Theorem}[section]
\newtheorem{proposition}[theorem]{Proposition}
\newtheorem{lemma}[theorem]{Lemma}
\newtheorem{corollary}[theorem]{Corollary}

\theoremstyle{definition}
\newtheorem{definition}[theorem]{Definition}
\newtheorem{example}[theorem]{Example}

\theoremstyle{remark}
\newtheorem{remark}[theorem]{Remark}

% ============================================================
% Notation
% ============================================================
\newcommand{\R}{\mathbb{R}}
\newcommand{\Z}{\mathbb{Z}}
\newcommand{\C}{\mathbb{C}}
\newcommand{\HH}{\mathbb{H}}
\newcommand{\OO}{\mathbb{O}}
\newcommand{\F}{\mathbb{F}}
\newcommand{\Stab}{\mathrm{Stab}}
\newcommand{\Hur}{\mathrm{Hur}}
\newcommand{\Cay}{\mathrm{Cay}}
\DeclareMathOperator{\rank}{rank}
\DeclareMathOperator{\Spin}{Spin}
\DeclareMathOperator{\Pin}{Pin}
\DeclareMathOperator{\SO}{SO}
\DeclareMathOperator{\SU}{SU}
\DeclareMathOperator{\U}{U}
\DeclareMathOperator{\SL}{SL}
\DeclareMathOperator{\Aut}{Aut}
\DeclareMathOperator{\Inn}{Inn}
\DeclareMathOperator{\Out}{Out}
\DeclareMathOperator{\Cl}{Cl}


% ============================================================
\title[Kinematic Spinors on Division-Algebra Root Systems]
{Kinematic Spinors on Division-Algebra Root Systems:\\
A First-Principles Derivation\\
of the Standard Model Weyl Group}

\author{Klee Irwin}
\address{Quantum Gravity Research, Los Angeles, CA}
\email{klee@quantumgravityresearch.org}

\date{\today}

\subjclass[2020]{17B22, 20F55, 51F25, 81R05}
\keywords{Root systems, Weyl groups, division algebras, Hopf fibrations,
Clifford algebras, spinors, Standard Model Weyl group}

% ============================================================
\begin{document}
% ============================================================

\begin{abstract}
We introduce \emph{kinematic spinors}: finite cyclic orbits of
$\Spin(n)$ rotors acting on root-system lattices, characterized
by a \emph{spinorial signature} in which the orbit in the
spinor module has twice as many elements as its image under
the $\Spin(n) \to \SO(n)$ double cover.
We show that kinematic spinors on the integral lattices of the
four normed division algebras
$\R, \C, \HH, \OO$ produce a compositional hierarchy of root
systems $A_1 \subset A_2 \subset D_4 \subset E_8$,
where each root system is a compound of oriented copies of
the previous one, generated by the spinor orbit at that level.

The 240 roots of $E_8$ decompose under the quaternionic Hopf
fibration $S^7 \to S^4$ into exactly ten copies of the $D_4$
root system, arranged in five mutually orthogonal pairs.
We prove computationally that this partition has an orbit of
$15{,}120$ elements under the Weyl group $W(E_8)$,
and that its stabilizer fits in a non-split extension
\[
  1 \;\longrightarrow\; 2T \;\longrightarrow\;
  \Stab_{W(E_8)}(\mathcal{P}_0) \;\longrightarrow\;
  W(D_5) \;\longrightarrow\; 1,
\]
where $W(D_5)\cong (\Z/2)^4 \rtimes S_5$ is the Weyl group of
$\SO(10)$ --- the canonical grand unified group --- and
$2T \cong \SL(2,\F_3)$ is the binary tetrahedral group.
The GUT breaking chain
$W(D_5) \supset W(A_4) \supset W(A_2)\times W(A_1)$,
with indices $16$ and $\binom{5}{3}=10$,
recovers the Standard Model Weyl group at each stage.

Triality --- the $S_3$ outer automorphism of $D_4$ that permutes
its vector, spinor, and conjugate-spinor representations ---
arises geometrically from the $4 = 1 + 3$ projection degeneracy
of the four hexagonal Hopf fibers of the $D_4$ root system.
We prove that orthogonal projection of $D_4$ along a fiber
normal produces a cuboctahedral shell whose pairwise angles
are exactly the $D_4$ angle set $\{60^\circ, 90^\circ,
120^\circ, 180^\circ\}$, preserving the root-system angular
structure in three dimensions.

All constructions use only discrete lattice-preserving cyclic
actions --- no continuous parameters, no free choices.
The Standard Model Weyl group emerges from the rigid geometry
of these kinematic spinors or not at all.
\end{abstract}

\maketitle

% ============================================================
\section{Introduction}\label{sec:intro}
% ============================================================

\subsection{The problem}

The Standard Model of particle physics is built on the gauge group
$\SU(3)_C \times \SU(2)_L \times \U(1)_Y$,%
\footnote{Strictly, the Standard Model gauge group is
$(\SU(3) \times \SU(2) \times \U(1)_Y)/\Z_6$;
the quotient by this discrete center is invisible at the
Weyl-group level and does not affect our results.}
where the subscripts denote color, left-chirality, and hypercharge.
This group, together with its three-generation fermion content,
has been confirmed by experiment to extraordinary precision.
Yet the group itself remains unexplained:
why these factors, why these representations, why three generations?

Approaches to deriving the Standard Model from deeper structure
broadly fall into two categories.
In the \emph{continuous approach}, one embeds
$\SU(3)\times\SU(2)\times\U(1)$ in a larger Lie group ---
$\SU(5)$ \cite{GeorgiGlashow1974},
$\SO(10)$ \cite{FritzschMinkowski1975},
$E_6$, or $E_8$ \cite{Lisi2007} ---
and breaks the larger symmetry via a Higgs mechanism.
The difficulty is that continuous Lie groups offer vast freedom:
with enough parameters, one can engineer nearly any subgroup structure.

In the \emph{algebraic approach}
\cite{Furey2016,Dixon1994,Boyle2020},
one works with the normed division algebras
$\R,\C,\HH,\OO$ and seeks the Standard Model in their
algebraic structure. This approach is more constrained but has not
produced a complete derivation linking all three gauge factors and
three generations to a single construction.

\subsection{Prior approaches}\label{sec:priorapproaches}

The problem of deriving the Standard Model gauge group from
deeper mathematical structure has motivated several distinct
research programs, which we briefly survey to situate the
present work.

\medskip
\noindent\textbf{Division-algebra representations.}\quad
Dixon \cite{Dixon1994} and Furey \cite{Furey2016} observe that
the tensor product $\C\otimes\HH\otimes\OO$ of the three
associative-plus-one normed division algebras contains algebraic
structures reproducing the quantum numbers of a single generation
of Standard Model fermions.
This program shares our algebraic starting point but works at the
level of representation content within a fixed algebraic substrate,
constructing the three gauge factors from separate algebraic
operations rather than from a unified geometric mechanism.

\medskip
\noindent\textbf{Exceptional Lie-algebra embedding.}\quad
Lisi \cite{Lisi2007} proposed embedding all Standard Model fields,
together with gravity, in the $E_8$ Lie algebra.
Distler and Garibaldi \cite{DistlerGaribaldi2009} showed that
the required chiral fermion spectrum cannot be accommodated within
$E_8$ representations.
Our construction avoids this obstruction entirely: we make no
claims about $E_8$ representations or continuous embeddings,
working exclusively with the discrete Weyl group $W(E_8)$ and
its action on lattice partitions.

\medskip
\noindent\textbf{Jordan-algebraic and characterization approaches.}\quad
Todorov and Dubois-Violette \cite{TodorovDuboisViolette2018}
deduce the Standard Model gauge symmetry from the automorphism
and structure groups of the exceptional Jordan algebra
$J_3(\OO)$.
Boyle \cite{Boyle2020} connects the same algebra to the
Standard Model with triality playing a central role.
Krasnov \cite{Krasnov2019} characterizes the Standard Model
gauge group as the subgroup of $\SO(9)$ preserving a certain
algebraic structure.
These approaches identify elegant algebraic contexts in
which the gauge group naturally lives, but each requires
selecting a particular subalgebra, projection, or
characterizing condition within the ambient structure.

\medskip
\noindent\textbf{Clifford spinor induction.}\quad
Closest to our framework is Dechant \cite{Dechant2016}, who
showed that ADE root systems arise as orbits of seed vectors
under finite subgroups of $\Spin(n)$ --- a ``spinor induction''
mechanism.
In particular, Dechant constructs $E_8$ from the icosahedral
group in $\Spin(4)$.
Our kinematic spinor approach shares this Clifford-algebraic
starting point but extends it in two directions:
(i)~we build a \emph{recursive}, multi-level hierarchy
($A_1 \to A_2 \to D_4 \to E_8$) rather than a one-step
generation, and
(ii)~we extract the Standard Model Weyl group from the Hopf
partition stabilizer at the top level, connecting discrete
lattice geometry to gauge-group structure.

\medskip
\noindent\textbf{The present contribution.}\quad
We show that the division-algebra root-system hierarchy,
analyzed through its discrete lattice-preserving symmetries
alone, forces the Standard Model Weyl group
$W(A_2)\times W(A_1) = S_3 \times S_2$
with no continuous parameters and no free choices.
By the classical reconstruction theorem for simply-laced root
systems (Remark~\ref{rem:reconstruction} in
\S\ref{sec:discussion}),
this discrete Weyl-group derivation canonically recovers the
continuous semisimple factor $\SU(3)\times\SU(2)$, narrowing
the gap to only the abelian factor $\U(1)_Y$ and the global
$\Z_6$ quotient.

\subsection{Discreteness as structural information}
\label{sec:discreteness}

A key insight motivating our approach is that continuous
symmetry groups, despite their mathematical elegance,
carry \emph{no internal structure}.
A continuous sphere $S^n$ is a homogeneous space:
every point is equivalent to every other point under $\SO(n+1)$.
There is nothing to distinguish, nothing to count, nothing
to decompose.

A root system, by contrast, is a finite set of points on a
sphere, constrained by integrality and crystallographic
conditions.
The finiteness breaks the continuous symmetry and introduces
structure: angles between roots, nesting of sub-root-systems,
partition counts, stabilizer groups.
Each distinct configuration carries specific combinatorial and
group-theoretic information that continuous objects cannot.

This observation has a precise consequence.
A continuous Lie group $G$ contains uncountably many
one-parameter subgroups, and embedding a target subgroup
$H \hookrightarrow G$ requires choosing continuous parameters
(breaking directions, vacuum expectation values).
The ``derivation'' of $H$ from $G$ is only as compelling as
the justification for these choices.

By contrast, the root system $E_8$ is a rigid combinatorial
object: $240$ specific vectors in $\R^8$ with prescribed inner
products.
There is no freedom in the object itself.
When we partition $E_8$ into $10$ copies of $D_4$ and compute
the stabilizer, the result $W(D_5) = W(\SO(10))$ is forced
by the geometry --- no parameters are chosen.
The subsequent breaking to the Standard Model Weyl group
$W(A_2) \times W(A_1)$ proceeds through a finite chain of
discrete inclusions, each with a determined index.

The restriction to discrete lattice-preserving actions is severe
but deliberate:
it is precisely this severity that makes the results compelling.
Either the lattice geometry produces the Standard Model structure
or it does not.
It produces the Standard Model structure.


\subsection{Main results}

We establish the following chain of results, each verified by
exact machine computation on the integral lattices.
The connecting thread is the kinematic spinor
(Definition~\ref{def:kinematic}):
a finite cyclic orbit in a spinor module that carries the
double-cover signature of the spin group.

\begin{theorem}[Compositional Hierarchy]\label{thm:A}
The root systems of the four normed division algebras nest
compositionally:
\begin{enumerate}[label=(\alph*)]
\item $A_2$ is a $3$-fold composite of $A_1$ (unique partition),
  generated by a kinematic spinor in $\Spin(2)$.
\item $D_4$ is a $4$-fold composite of $A_2$ (exactly $8$
  partitions), organized by the complex Hopf fibration
  $S^3 \to S^2$.
\item $E_8$ is a $10$-fold composite of $D_4$ (exactly $15{,}120$
  partitions under $W(E_8)$), organized by the quaternionic
  Hopf fibration $S^7 \to S^4$.
\end{enumerate}
\end{theorem}

\begin{theorem}[Partition Stabilizer]\label{thm:B}
The stabilizer of the $E_8$ Hopf partition in $W(E_8)$ has order
$46{,}080$ and fits in a non-split exact sequence:
\[
  1 \to 2T \to \Stab(\mathcal{P}_0) \to W(D_5) \to 1,
\]
where $W(D_5)$ is the Weyl group of $\SO(10)$ and
$2T \cong \SL(2,\F_3)$ is the binary tetrahedral group of
order~$24$.
The kernel $2T$ is the group of unit Hurwitz quaternions ---
a finite subgroup of $\Spin(3)$ --- acting as a kinematic spinor
on each $D_4$ fiber.
\end{theorem}

\begin{theorem}[GUT Chain]\label{thm:C}
The Weyl-group-level GUT breaking chain
\[
  W(D_5) \supset W(A_4) \supset W(A_2)\times W(A_1)
\]
with indices $16$ and $10$ recovers the Weyl groups of
$\SO(10) \supset \SU(5) \supset \SU(3)\times\SU(2)\times\U(1)$.
\end{theorem}

\begin{theorem}[Triality from Projection]\label{thm:D}
The four $A_2$ Hopf fibers of $D_4$ lie in four
$2$-dimensional subspaces whose normal directions form a
regular tetrahedron on $S^3$.
Orthogonal projection along any fiber normal breaks
$S_4 \to 1 + S_3$, where the residual $S_3$ coincides with
$\Out(D_4) \cong S_3$.
\end{theorem}

\begin{theorem}[Semi-Conformal Angle Preservation]
\label{thm:E}
Orthogonal projection of $D_4$ along a Hopf fiber normal
separates the $24$ roots into two radial shells:
\begin{enumerate}[label=(\alph*)]
\item An outer shell of $12$ roots forming a
  cuboctahedron (root system $A_3$), with pairwise angles
  exactly $\{60^\circ, 90^\circ, 120^\circ, 180^\circ\}$
  --- the full $D_4$ angle set.
\item An inner shell of $12$ roots projecting to $6$ distinct
  positions forming an octahedron, with pairwise angles
  $\{90^\circ, 180^\circ\}$.
\end{enumerate}
\end{theorem}

\subsection{Outline}

Section~\ref{sec:spinor} defines kinematic spinors and
establishes the double-cover signature.
Section~\ref{sec:hierarchy} proves the division-algebra
compositional hierarchy.
Section~\ref{sec:hopf} constructs the quaternionic Hopf
partition of $E_8$.
Section~\ref{sec:stabilizer} computes the partition count
and stabilizer.
Section~\ref{sec:gut} extracts the GUT chain.
Section~\ref{sec:triality} derives triality and proves the
semi-conformal projection theorem.
Section~\ref{sec:discussion} discusses implications and open
questions.


% ============================================================
\section{Kinematic Spinors on Root Lattices}
\label{sec:spinor}
% ============================================================

\subsection{Clifford algebras and spin groups}

We briefly recall the algebraic setting.
The \emph{Clifford algebra} $\Cl(n)$ is the associative
algebra generated by $\R^n$ subject to the relation
$v \cdot v = -\|v\|^2$ for all $v \in \R^n$
\cite{LawsonMichelsohn1989}.
For an orthonormal basis $\{e_1,\ldots,e_n\}$, this gives
$e_i e_j + e_j e_i = -2\delta_{ij}$.
The algebra $\Cl(n)$ has dimension $2^n$ and contains
$\R^n$ as a linear subspace.

The \emph{spin group} $\Spin(n) \subset \Cl(n)$ is the
subgroup of even-grade invertible elements $g$ satisfying
$g v g^{-1} \in \R^n$ for all $v \in \R^n$ and
$g \tilde{g} = 1$ (where $\tilde{g}$ denotes the reversal
anti-automorphism).
There is an exact sequence
\begin{equation}\label{eq:doublecover}
  1 \;\longrightarrow\; \Z/2 \;\longrightarrow\;
  \Spin(n) \;\xrightarrow{\;\rho\;}\;
  \SO(n) \;\longrightarrow\; 1,
\end{equation}
where the map $\rho$ sends $g$ to the orthogonal transformation
$v \mapsto g v g^{-1}$, and
$\ker(\rho) = \{+1, -1\} \subset Z(\Spin(n))$.
(The full center of $\Spin(n)$ depends on $n \bmod 4$
and may be larger; only the central $\Z/2 = \ker(\rho)$
is used in our arguments.)

A \emph{spinor module} $S$ for $\Cl(n)$ is an irreducible
left module.
Every element $g \in \Spin(n) \subset \Cl(n)$ acts on $S$ by
left multiplication: $\psi \mapsto g \cdot \psi$.
This is the \emph{spinor representation} of $\Spin(n)$,
which does not factor through $\SO(n)$: the kernel element
$-1 \in \Spin(n)$ acts as $\psi \mapsto -\psi$, not as the
identity.

\subsection{Kinematic spinors: definition}

The central object of this paper is a discrete analog of the
spinor orbit.

\begin{definition}[Kinematic spinor]\label{def:kinematic}
Let $g \in \Spin(n)$ be an element of finite order $m$
(so $g^m = 1$ in $\Spin(n)$),
and let $S$ be a spinor module for $\Cl(n)$.
For a nonzero $\psi \in S$ whose stabilizer in
$\langle g \rangle$ is trivial
(i.e.\ $g^k \psi \neq \psi$ for $0 < k < m$),
the \emph{kinematic spinor}
generated by $(g, \psi)$ is the finite cyclic orbit
\[
  \mathcal{O}_S(g,\psi) \;=\;
  \bigl\{g^k \cdot \psi : k = 0, 1, \ldots, m-1\bigr\}
  \;\subset\; S.
\]
This orbit has exactly $m$ elements.
We say $g$ has \emph{spinorial order} $m$ if $g^m = 1$
in $\Spin(n)$ and $g^j \neq 1$ for $0 < j < m$.
\end{definition}

The qualifier \emph{kinematic} reflects the fact that the full
orbit --- not any individual point --- carries the
group-theoretic information.
Just as a kinematic trajectory in mechanics is determined by
the sequence of positions traversed under a dynamical rule,
a kinematic spinor is the discrete trajectory of a spinor
under repeated application of a Clifford rotor.
The configuration $\mathcal{O}_S(g,\psi)$ encodes the
structure of the cyclic group
$\langle g \rangle \subset \Spin(n)$, including its
spinorial character.
The construction is related to but distinct from Dechant's
\emph{spinor induction} \cite{Dechant2016}, which generates
a root system as a single orbit of a finite subgroup of
$\Spin(n)$.
Our construction is recursive (each level's output seeds
the next) and focuses on the orbit structure and partition
stabilizers rather than root-system generation per~se.

\begin{remark}\label{rem:groupkinematic}
Definition~\ref{def:kinematic} uses a single element
$g \in \Spin(n)$ generating a cyclic orbit.
More generally, for a finite subgroup
$\Gamma \subset \Spin(n)$ and $\psi \in S$, the
\emph{$\Gamma$-kinematic spinor} is the orbit
$\Gamma \cdot \psi = \{\gamma \cdot \psi : \gamma \in \Gamma\}
\subset S$.
This generalization is needed at Level~3 of the hierarchy,
where the kernel $2T \cong \SL(2,\F_3)$ --- a non-cyclic
finite subgroup of $\Spin(3)$ --- acts on each $D_4$ fiber.
Each cyclic subgroup $\langle \gamma \rangle \subset 2T$
generates a kinematic spinor in the sense of
Definition~\ref{def:kinematic}; the full $2T$-orbit is
their union.
\end{remark}

\begin{remark}
The definition is elementary: $\mathcal{O}_S(g,\psi)$ is
the orbit of $\psi$ under a cyclic group action on a module.
What is non-trivial is the interplay between the spinorial
order $m$ and the \emph{vector order} of $g$'s image in
$\SO(n)$, which we now characterize.
\end{remark}

\subsection{The double-cover signature}

\begin{proposition}[Spinorial signature]\label{prop:signature}
Let $g \in \Spin(n)$ have even spinorial order $m = 2p$,
and suppose $g^p = -1$ in $\Spin(n)$
(i.e., the half-order element is the nontrivial kernel element
of the double cover).
Let $\bar{g} = \rho(g) \in \SO(n)$ be the image of $g$.
Then:
\begin{enumerate}[label=(\roman*)]
\item $\bar{g}$ has order exactly $p = m/2$ in $\SO(n)$.
\item For any $\psi \in S$ whose stabilizer in
  $\langle g \rangle$ is trivial
  (i.e.\ $g^k \psi \neq \psi$ for $0 < k < m$),
  the spinor orbit
  $\mathcal{O}_S(g,\psi)$ has exactly $m$ elements.
\item For any vector $v \in \R^n$ whose stabilizer in
  $\langle \bar{g} \rangle$ is trivial
  (i.e.\ $\bar{g}^k v \neq v$ for $0 < k < p$),
  the vector orbit
  $\mathcal{O}_V(\bar{g},v) =
  \{\bar{g}^k v : k = 0,\ldots,p-1\}$
  has exactly $p = m/2$ elements.
\end{enumerate}
For such generic pairs $(\psi, v)$, the ratio
\begin{equation}\label{eq:signature}
  \frac{|\mathcal{O}_S(g,\psi)|}{|\mathcal{O}_V(\bar{g},v)|}
  \;=\; 2
\end{equation}
is the \emph{spinorial signature}: the factor of~$2$ that
distinguishes spinor orbits from vector orbits, directly
reflecting the double cover~\eqref{eq:doublecover}.
\end{proposition}

\begin{proof}
(i) Since $g^p = -1$, we have
$\bar{g}^p = \rho(-1) = I$, so
$\mathrm{ord}(\bar{g})$ divides $p$.
Suppose $\mathrm{ord}(\bar{g}) = d$ with $d \mid p$
and $d < p$.
Then $g^d \in \ker\rho = \{+1,-1\}$.
If $g^d = +1$, then $\mathrm{ord}(g) \mid d < 2p = m$,
contradicting $\mathrm{ord}(g) = m$.
If $g^d = -1$, then $g^{2d} = 1$, so
$\mathrm{ord}(g) \mid 2d$; but $2d < 2p = m$,
again contradicting $\mathrm{ord}(g) = m$.
Hence $\mathrm{ord}(\bar{g}) = p$.

(ii) By hypothesis,
$\Stab_{\langle g\rangle}(\psi) = \{1\}$,
so the orbit--stabilizer theorem gives
$|\mathcal{O}_S| = |\langle g\rangle| = m$.
(The genericity condition excludes $\psi$ lying in
a proper eigenspace of some $g^k$; for each
$0 < k < m$, the fixed locus of $g^k$ on $S$ is a
proper subspace, so generic $\psi$ satisfies the hypothesis.)

(iii) By hypothesis,
$\Stab_{\langle\bar{g}\rangle}(v) = \{1\}$,
so $|\mathcal{O}_V| = |\langle\bar{g}\rangle| = p$.
(Note that $v$ not fixed by $\bar{g}$ is necessary but
not sufficient: $v$ must also avoid the fixed subspaces of
$\bar{g}^k$ for all $0 < k < p$.)

The ratio is $m/p = 2$.
\end{proof}

\begin{remark}[Genericity is automatic in all instances]
\label{rem:genericity}
In each concrete instance used in this paper, the
stabilizer-trivial condition is satisfied by every
nonzero seed, for the following reasons:
\begin{itemize}
\item \textbf{Level~1} ($g = \exp(\tfrac{\pi}{3} e_1 e_2)
  \in \Spin(2)$, order~$6$):
  The spinor module is $S \cong \R^2$, and $g$ acts as
  rotation by $\pi/3$.
  The only fixed point of a nontrivial rotation on $\R^2$
  is the origin, so every nonzero $\psi$ has trivial
  stabilizer.
\item \textbf{Level~2} ($\U(1) \subset \Spin(4)$
  fiber action):
  Each $\U(1)$ subgroup acts on a $2$-plane by rotation
  of angle $2\pi/3$ (a primitive $3$rd root of unity).
  A nonzero vector in this plane is fixed by the $k$-th
  power only if $k \equiv 0 \pmod{3}$, so the stabilizer
  in $\langle g \rangle$ is trivial for every nonzero
  root vector in the fiber.
\item \textbf{Level~3} ($2T \cong \SL(2,\F_3) \subset
  \Spin(3) \cong \HH^1$, order~$24$):
  Each element $q \in 2T$ acts on $\HH \cong \R^4$ by
  left multiplication $v \mapsto q \cdot v$.
  If $q \cdot v = v$ for some $v \neq 0$, then
  $(q - 1) \cdot v = 0$; since $\HH$ is a division
  algebra (no zero divisors) and $v \neq 0$, this
  forces $q - 1 = 0$, i.e.\ $q = 1$.
  Hence every nontrivial element of $2T$ acts
  fixed-point-freely on $\HH \setminus \{0\}$.
\end{itemize}
Thus, in all three levels, the genericity conditions of
(ii) and (iii) are automatically satisfied by
\emph{any} nonzero seed.
\end{remark}

\begin{remark}
If $g$ has odd order $m$ in $\Spin(n)$, then $g^k = -1$
is impossible for any $k$ (since $(-1)^2 = 1$ would give
$2k \equiv 0 \pmod{m}$, contradicting $m$ odd and $0 < 2k < 2m$).
In this case $\bar{g}$ also has order $m$, and the spinor
and vector orbits have the same cardinality: there is no
spinorial signature.
The kinematic spinors of interest in this paper all have
even order with the spinorial condition $g^{m/2} = -1$.
\end{remark}

\begin{remark}
The spinorial signature~\eqref{eq:signature} is the discrete
manifestation of the defining property of spinors: they
transform under the double cover of the rotation group.
In the continuous setting, a $2\pi$ rotation acts as $-1$ on
spinors and $+1$ on vectors.
In our discrete setting, the half-order element $g^{m/2} = -1$
acts as $-1$ on spinors and $+1$ on vectors, producing
exactly twice as many distinct spinor orbit points as vector
orbit points.
\end{remark}

\subsection{Example: Level 1 spinor}
\label{sec:level1}

\begin{example}\label{ex:level1}
Let $R = \exp(\tfrac{\pi}{3}\, e_1 e_2) \in \Spin(2)$.
In the Clifford algebra $\Cl(2)$, the bivector $e_1 e_2$
satisfies $(e_1 e_2)^2 = -1$, so
\[
  R \;=\; \cos\tfrac{\pi}{3} + \sin\tfrac{\pi}{3}\; e_1 e_2
    \;=\; \tfrac{1}{2} + \tfrac{\sqrt{3}}{2}\, e_1 e_2.
\]
The spinorial order of $R$ is $6$:
$R^6 = \exp(2\pi\, e_1 e_2) = 1$, and
$R^3 = \exp(\pi\, e_1 e_2) = -1$ (the kernel element of
$\Spin(2) \to \SO(2)$).

\emph{Spinor orbit.}
The spinor module $S$ for $\Cl(2)$ is $\R^2$, with
$e_1 e_2$ acting as the matrix
$J = \bigl(\begin{smallmatrix} 0 & -1 \\ 1 & 0
\end{smallmatrix}\bigr)$.
Taking $\psi = (1, 0)$, the orbit
$\mathcal{O}_S(R, \psi) = \{R^k \cdot \psi : k=0,\ldots,5\}$
is a regular hexagon with $6$ vertices --- after rescaling
to squared norm~$2$, this is the root system $A_2$.

\emph{Vector orbit.}
The image $\bar{R} = \rho(R) \in \SO(2)$ is rotation by
$2\pi/3$.
The orbit of any vector under $\bar{R}$ has $3$ elements
(an equilateral triangle), which after including negatives
gives $A_2$ from the vector perspective.

The spinorial signature is $6/3 = 2$:
the \emph{kinematic spinor orbit is the full $A_2$ root system},
while the vector orbit gives only half of it.
This is Level~1 of the hierarchy: the root system $A_2$
is literally the kinematic spinor of a $\Spin(2)$ rotor
acting on $A_1$ seeds.
\end{example}


% ============================================================
\section{The Division-Algebra Root System Hierarchy}
\label{sec:hierarchy}
% ============================================================

\subsection{Division algebras and their integral forms}

The four normed division algebras over $\R$ ---
the reals $\R$, complexes $\C$, quaternions $\HH$, and
octonions $\OO$ ---
have dimensions $1, 2, 4, 8$ respectively.
By the theorems of Hurwitz \cite{Hurwitz1898} (composition algebras)
and Adams \cite{Adams1960} (Hopf invariant one),
these are the only normed division algebras and the only dimensions
admitting Hopf fibrations.

Each normed division algebra admits multiple integral forms
(algebraic rings of integers), and the choice matters.
For $\C$, one may take the Gaussian integers $\Z[i]$
(whose $4$ units form the root system $B_1 \times B_1$)
or the Eisenstein integers $\Z[\omega]$
(whose $6$ units form $A_2$).
For $\HH$, one may take the Lipschitz integers
$\Z\langle 1,i,j,k\rangle$ (whose $8$ units form $D_2 \times D_2$),
or the Hurwitz integers $\Hur(\HH)$ (whose $24$ units form $D_4$).
For $\OO$, several non-isomorphic maximal orders exist,
all related by $\Aut(\OO) \cong G_2$, but their unit groups
are isomorphic.

In each case, we select the integral form whose unit lattice
is the \emph{densest lattice sphere packing} in the
corresponding dimension.
This is a canonical choice: the densest lattice sphere packing
in $\R^n$ is unique up to isometry (among lattice packings)
for $n = 1, 2, 4, 8$
\cite{ConwaySloane1999,Viazovska2017}.
The norm-one units of this maximal-density integral form
generate a root system
(after the standard rescaling to squared norm~$2$):

\medskip
\begin{center}
\begin{tabular}{@{}lllccl@{}}
\toprule
Algebra & Integral form & Dim & Root system & $|\text{Roots}|$
  & Units \\
\midrule
$\R$ & $\Z$ & 1 & $A_1$ & 2 & 2 \\
$\C$ & $\Z[\omega]$, $\omega = e^{2\pi i/3}$ & 2 & $A_2$ & 6
  & 6 \\
$\HH$ & $\Hur(\HH) = \Z\langle 1,i,j,\tfrac{1+i+j+k}{2}\rangle$
       & 4 & $D_4$ & 24 & 24 \\
$\OO$ & $\Cay(\OO)$ & 8 & $E_8$ & 240 & 240 \\
\bottomrule
\end{tabular}
\end{center}
\medskip

The Eisenstein integers $\Z[\omega]$ are the ring of integers in
$\mathbb{Q}(\sqrt{-3})$; their $6$ units are the
$6$th roots of unity.
The Hurwitz quaternions $\Hur(\HH)$ are the maximal order in the
rational quaternion algebra $\HH(\mathbb{Q})$;
their 24 units are the vertices of the 24-cell
\cite{ConwaySloane1999}.
The Cayley integers $\Cay(\OO)$ form the (unique up to
$\Aut(\OO)$) maximal order in the rational octonion algebra;
their 240 units are the $E_8$ root vectors
\cite{ConwaySloane1999,Coxeter1946}.

\begin{remark}[The maximal-density selection principle]
\label{rem:packing}
The unit lattices $A_1, A_2, D_4, E_8$ are precisely the
kissing configurations of the densest lattice sphere packings
in $\R^1, \R^2, \R^4, \R^8$ respectively
\cite{ConwaySloane1999,Viazovska2017}.
This is not a coincidence: the norm-multiplicativity of
division algebras forces the maximal-density integral form
to have the largest possible number of units.
Among all integral forms of a given division algebra,
the one whose unit lattice is the densest packing has the
most unit roots and therefore the richest algebraic
structure.

The choice is critical: replacing the Eisenstein integers
with the Gaussian integers would give $A_1 \times A_1$
instead of $A_2$ --- losing the hexagonal structure that
generates the hierarchy.
Replacing the Hurwitz integers with the Lipschitz integers
would give a root system with $8$ instead of $24$ units ---
losing the 24-cell geometry that produces the Hopf partition.
The compositional hierarchy we prove below exists
\emph{only} for the maximal-density chain
$\Z \subset \Z[\omega] \subset \Hur(\HH) \subset \Cay(\OO)$,
and is destroyed by any other choice of integral form.

From this perspective, the hierarchy is a statement about how
maximally dense sphere packings nest across the
division-algebra dimensions --- and the Standard Model
structure emerges because nature's geometry selects the
densest packing at each level, or equivalently,
the integral form with the maximum number of units.
\end{remark}

\subsection{Compositional partitions}

\begin{definition}\label{def:compositional}
Let $\Phi$ be a root system in $\R^n$ and $\Psi$ a root system
of rank $< n$.
A \emph{compositional partition of $\Phi$ into copies of $\Psi$}
is a decomposition
$\Phi = \Psi_1 \sqcup \Psi_2 \sqcup \cdots \sqcup \Psi_k$
(disjoint union)
such that each $\Psi_j$ is a root system in some subspace
$V_j \subset \R^n$ that is isometric (as a root system) to $\Psi$.
The integer $k = |\Phi|/|\Psi|$ is the \emph{composite index}.
We use the term ``compositional'' to emphasize that
$\Phi$ is assembled from oriented copies of $\Psi$
in distinct subspaces, in the manner of a geometric compound.
\end{definition}

\subsection{Level 1: $A_2 = 3 \times A_1$ via $\Spin(2)$}

\begin{proposition}[$A_2 = 3 \times A_1$]\label{prop:A2}
The root system $A_2$ admits a unique compositional partition into
$3$ copies of $A_1$.
\end{proposition}

\begin{proof}
The six roots of $A_2$ form three antipodal pairs
$\{r, -r\}$.
Each pair is a copy of $A_1$ in a one-dimensional subspace.
Since every root belongs to exactly one antipodal pair,
the partition is unique.
\end{proof}

\begin{remark}[Spinor interpretation of Level~1]
\label{rem:spinorlevel1}
By Example~\ref{ex:level1}, this partition is generated by the
kinematic spinor of $R = \exp(\tfrac{\pi}{3}\, e_1 e_2)
\in \Spin(2)$.
The rotor $R$ has spinorial order $6$ and $R^3 = -1$.
Its orbit carries a single $A_1$ seed through
$3$ orientations at $60^\circ$ spacing, assembling the full
$A_2$ hexagon.
The spinorial signature $6/3 = 2$ reflects the fact that
the spinor representation of $\Spin(2)$ distinguishes $A_2$
(the full hexagon) from the vector representation's triangle.
This is the simplest instance of the paper's central mechanism:
a discrete spinor orbit on a root lattice generating the next
root system in the hierarchy.
\end{remark}

\subsection{Level 2: $D_4 = 4 \times A_2$ via $S^3 \to S^2$}

\begin{proposition}[$D_4 = 4 \times A_2$]\label{prop:D4}
The root system $D_4$ admits exactly $8$ compositional partitions
into $4$ copies of $A_2$.
Under the $S_3$ outer automorphism group $\Out(D_4)$, these $8$
partitions form two orbits: one of size $2$ and one of size $6$.
\end{proposition}

\begin{proof}
Machine-verified by exhaustive enumeration.

\emph{Step 1.}
We enumerate all $A_2$ sub-root-systems inside $D_4$
by identifying $6$-element subsets of the $24$ roots
satisfying the $A_2$ inner-product conditions and
spanning a $2$-dimensional subspace.
This yields $16$ distinct $A_2$ sub-root-systems.

\emph{Step 2.}
Exhaustive search for sets of $4$ mutually disjoint $A_2$'s
covering all $24$ roots yields exactly $8$ partitions.

\emph{Step 3.}
The $S_3$ triality group acts on these $8$ partitions,
producing two orbits: $\{P_0, P_1\}$ (each fixed by the
order-$3$ element $\tau$, swapped by $\sigma$) and
$\{P_2,\ldots,P_7\}$ (full $S_3$ orbit decomposing into
two interleaved $3$-cycles).
See Appendix~\ref{app:computation}.
\end{proof}

\begin{remark}[Spinor interpretation of Level~2]
\label{rem:spinorlevel2}
The $24$ roots of $D_4$ lie on $S^3(\sqrt{2}) \subset \R^4$.
The complex Hopf fibration $S^3 \to S^2$ partitions these
$24$ roots into the $4$ hexagonal $A_2$ fibers of
Proposition~\ref{prop:D4}.
Each fiber is a great circle on $S^3$, and the
$4$ fiber planes form a regular tetrahedron
(Proposition~\ref{prop:tetrahedron}).

The Hopf fibration is the geometric manifestation of a
$\U(1) \subset \Spin(4)$ action: each fiber is the orbit
of a root under the $\U(1)$ subgroup corresponding to
the fiber direction.
The $4$ choices of fiber direction correspond to the $4$
vertices of the tetrahedron, and the $8$ partitions arise
from the $4 \times 2$ choices (each fiber direction with
two orientations).
This is the Level~2 kinematic spinor:
the $\Spin(4)$ fiber structure organizes $D_4$ as a compound
of $4$ oriented $A_2$ copies.
\end{remark}

\subsection{Level 3: $E_8 = 10 \times D_4$ via $S^7 \to S^4$}

\begin{proposition}[$E_8 = 10 \times D_4$]\label{prop:E8}
The root system $E_8$ admits compositional partitions into $10$
copies of $D_4$.
Under $W(E_8)$, the orbit of the quaternionic Hopf partition
contains exactly $15{,}120$ elements.
\end{proposition}

\begin{remark}
The decomposition of $E_8$ into $10$ copies of $D_4$ via the
quaternionic Hopf fibration was first constructed in
\cite{ClawsonFangIrwin2024}; see also \cite{FangIrwin2024}.
The present paper extends this by computing the full partition
orbit, its stabilizer structure, the GUT breaking chain, and
the semi-conformal projection.
\end{remark}

\begin{proof}
Existence follows from the Hopf construction
(Section~\ref{sec:hopf}).
The orbit count is proved in Theorem~\ref{thm:orbit}.
\end{proof}

\begin{remark}[Spinor interpretation of Level~3]
\label{rem:spinorlevel3}
The quaternionic Hopf fibration $S^7 \to S^4$ is the
$\Spin(8)$-level analog: the fiber is $S^3$ (the group
manifold of $\Spin(3) \cong \SU(2)$), and each fiber
intersects the $E_8$ root system in a $D_4$ copy.
The $10$ fibers over the $10$ cross-polytope vertices
are the $10$ $D_4$ orientations.
The kernel of the partition stabilizer --- the binary
tetrahedral group $2T \cong \SL(2,\F_3)$ --- is a
finite subgroup of $\Spin(3)$ that acts as a kinematic spinor
on each $D_4$ shell (Theorem~\ref{thm:stab}).
\end{remark}


\subsection{Partition counts and the spinor hierarchy}

\begin{center}
\begin{tabular}{@{}llccll@{}}
\toprule
Level & Partition & Count & $|\Stab|$ & Fiber structure & Spinor \\
\midrule
1 & $A_2 = 3 \times A_1$ & 1 & 6 & $\Z/2$ (antipodal) & $\Spin(2)$, $R^3 = -1$ \\
2 & $D_4 = 4 \times A_2$ & 8 & 24 & $S^1 \hookrightarrow S^3 \to S^2$ & $\U(1) \subset \Spin(4)$ \\
3 & $E_8 = 10 \times D_4$ & 15{,}120 & 46{,}080 & $S^3 \hookrightarrow S^7 \to S^4$ & $2T \subset \Spin(3)$ \\
\bottomrule
\end{tabular}
\end{center}

\begin{remark}\label{rem:threehopf}
Levels~2 and~3 correspond to the complex and quaternionic
Hopf fibrations respectively.
Level~1 uses only the antipodal structure ($\Z/2$ acting
on $S^0 = \{+1, -1\}$), which is the fiber of the real
Hopf map $S^1 \to S^1$ but is simpler than a full Hopf
fibration.
Adams' theorem \cite{Adams1960} shows that Hopf fibrations
$S^{2n-1} \to S^n$ exist only for $n = 1, 2, 4, 8$
(the last being the octonionic fibration
$S^7 \hookrightarrow S^{15} \to S^8$).
The hierarchy terminates at $E_8$ not because the octonionic
Hopf fibration fails to exist, but because there is no
$16$-dimensional normed division algebra and hence no root
system beyond $E_8$ in the chain to decompose.
\end{remark}


% ============================================================
\section{The Quaternionic Hopf Partition of $E_8$}
\label{sec:hopf}
% ============================================================

\subsection{The Hopf fibration $S^7 \to S^4$}

The quaternionic Hopf map
$h \colon S^7 \to S^4$ is defined by
\cite{Baez2002,Hopf1935}:
\begin{equation}\label{eq:hopf}
  h(q_1, q_2) \;=\;
  \bigl(|q_1|^2 - |q_2|^2,\; 2\,q_1 \bar{q}_2\bigr)
  \;\in\; \R \times \R^4 \;=\; \R^5,
\end{equation}
where $q_1 = (x_0,x_1,x_2,x_3)$ and $q_2 = (x_4,x_5,x_6,x_7)$
are treated as quaternions, and $\bar{q}_2$ denotes quaternionic
conjugation.
The fiber $h^{-1}(p)$ over each point $p \in S^4$ is
diffeomorphic to $S^3$.

Since all $240$ roots of $E_8$ have squared norm $2$,
they lie on $S^7(\sqrt{2})$.
The Hopf map sends them to $S^4(2)$.

\begin{definition}\label{def:hopfpartition}
The \emph{$E_8$ Hopf partition} $\mathcal{P}_0$ is the partition
of the $240$ $E_8$ roots into fibers of the quaternionic
Hopf map~\eqref{eq:hopf}:
two roots $r, r'$ belong to the same cluster if and only if
$h(r) = h(r')$.
As we shall verify, this is a compositional partition
(Definition~\ref{def:compositional}) with $\Psi = D_4$
and composite index $k = 10$.
\end{definition}

\begin{proposition}[Hopf images]\label{prop:crosspolytope}
The Hopf images of the $240$ $E_8$ roots are exactly the
$10$ vertices of the cross-polytope $\beta_5$ in $\R^5$:
\[
  \{\pm 2\, e_i : i = 1,\ldots,5\}.
\]
Each vertex has exactly $24$ preimages among the $240$ roots.
\end{proposition}

\begin{proof}
Machine-verified.
We compute $h(r)$ for all $240$ roots $r$ of $E_8$
in the standard coordinate representation
($112$ integer-type roots $\pm e_i \pm e_j$,
$i < j$, plus $128$ half-integer-type roots
$(\pm\tfrac{1}{2},\ldots,\pm\tfrac{1}{2})$ with an even
number of minus signs).
All Hopf images have norm $2$ and cluster into exactly $10$
distinct points: the $5$ antipodal pairs
$\{+2 e_i, -2 e_i\}$, $i=1,\ldots,5$.
Each cluster contains exactly $24$ roots.
\end{proof}

\subsection{Each Hopf cluster is a $D_4$ root system}

\begin{proposition}\label{prop:D4clusters}
Each of the $10$ Hopf clusters is a root system of type $D_4$:
$24$ vectors spanning a $4$-dimensional subspace of $\R^8$,
with Cartan matrix of type $D_4$.
\end{proposition}

\begin{proof}
Machine-verified for each cluster $C_k$ ($k=1,\ldots,10$):
(i) the $24$ vectors span a $4$-dimensional subspace
(rank $= 4$ by exact row-reduction over $\mathbb{Q}$);
(ii) all pairwise inner products lie in $\{0, \pm 1, \pm 2\}$;
(iii) the set is closed under reflections through its own
elements (the root system axiom); and
(iv) the Cartan matrix of a simple system matches type $D_4$.
By the classification of irreducible crystallographic root systems
\cite{Humphreys1990}, a rank-$4$ root system with $24$ roots
and Cartan matrix of type $D_4$ is uniquely $D_4$.
\end{proof}

\subsection{The $5+5$ bipartite structure}

\begin{proposition}[Perpendicular pairs]\label{prop:perppairs}
The $10$ Hopf clusters form $5$ perpendicular pairs:
for each $i$, clusters $C_{+i}$ and $C_{-i}$ span
orthogonal $4$-dimensional subspaces of $\R^8$.
\end{proposition}

\begin{proof}
Machine-verified.
The cross-Gram matrix $G_{kl} = V_k V_l^T$ vanishes for
exactly $5$ pairs, with
$\dim(\operatorname{span}C_{+i}) +
\dim(\operatorname{span}C_{-i}) = 8$.
\end{proof}

\begin{remark}[The kinematic spinor perspective]
\label{rem:bipartitekinematic}
The $5+5$ bipartite structure is the central geometric object
of this paper.
Each perpendicular pair spans all of $\R^8$:
the two $D_4$ root systems in a pair are ``complementary
views'' of the $8$-dimensional space from orthogonal
$4$-dimensional vantage points.
This is the Level~3 kinematic spinor structure:
the quaternionic Hopf fiber ($S^3$, the spinor space of
$\Spin(3)$) mediates between the two complementary $D_4$'s,
with the binary tetrahedral group $2T$ acting as the discrete
spinor symmetry on each fiber.
\end{remark}


% ============================================================
\section{The Partition Count and Stabilizer}
\label{sec:stabilizer}
% ============================================================

\subsection{The orbit under $W(E_8)$}

The Weyl group $W(E_8)$ is generated by reflections
\begin{equation}\label{eq:reflection}
  s_\alpha(r) = r - \langle r,\alpha\rangle\,\alpha
\end{equation}
for roots $\alpha$ of $E_8$ (using
$\langle\alpha,\alpha\rangle = 2$).
It has order $696{,}729{,}600 = 2^{14}\cdot 3^5\cdot 5^2\cdot 7$
\cite{Humphreys1990}.

\begin{theorem}[Orbit size]\label{thm:orbit}
Under $W(E_8)$, the Hopf partition $\mathcal{P}_0$ has an
orbit of exactly $15{,}120$ elements.
\end{theorem}

\begin{proof}
Machine-verified by breadth-first search (BFS).
Starting from $\mathcal{P}_0$, all $120$ root reflections
(one per antipodal pair) are applied iteratively.
Since $W(E_8)$ is generated by reflections, the BFS explores
the full orbit.
The queue drains after visiting $15{,}120$ distinct
partitions, proving orbit completeness.

By the orbit-stabilizer theorem,
$|\Stab(\mathcal{P}_0)| = 696{,}729{,}600 / 15{,}120
= 46{,}080$.
Combinatorially,
$15{,}120 = 9!/4! = P(9,5)$.
\end{proof}

\begin{remark}\label{rem:orderidentity}
The factorization
$|W(E_8)| = 192 \times 3{,}628{,}800 = |W(D_4)| \times |S_{10}|$
reflects the composite structure $E_8 = 10 \times D_4$:
the Weyl group is ``large enough'' for the $10!$-fold
permutations of orientations, modulo the fiber stabilizer.
This is an \emph{order identity}, not a group isomorphism ---
$W(E_8)$ is \emph{not} isomorphic to $W(D_4) \times S_{10}$
(the extension~\eqref{eq:ses} is non-split).
\end{remark}

\subsection{The stabilizer: structure theorem}

\begin{theorem}[Stabilizer structure]\label{thm:stab}
The stabilizer $\Stab(\mathcal{P}_0)$ fits in a non-split
short exact sequence
\begin{equation}\label{eq:ses}
  1 \;\longrightarrow\; K \;\longrightarrow\;
  \Stab(\mathcal{P}_0) \;\xrightarrow{\;\varphi\;}
  W(D_5) \;\longrightarrow\; 1,
\end{equation}
where:
\begin{enumerate}[label=(\roman*)]
\item $\varphi$ records how each stabilizer element permutes
  the $10$ clusters.
\item $\operatorname{im}\varphi \cong W(D_5)
  = (\Z/2)^4 \rtimes S_5$, order $1{,}920$.
\item $\ker\varphi \cong 2T \cong \SL(2,\F_3)$, order~$24$.
\item The extension is non-split:
  $K$ is not central in $\Stab(\mathcal{P}_0)$,
  and no complement $H \cong W(D_5)$ with $K \cap H = \{1\}$
  exists.
\end{enumerate}
\end{theorem}

\begin{proof}
Machine-verified via the Schreier generator method
\cite{Sims1970}.

\emph{Step 1: Construct the stabilizer.}
The BFS of Theorem~\ref{thm:orbit} records coset representatives
$g_j$ satisfying $g_j(\mathcal{P}_0) = \mathcal{P}_j$.
Schreier generators $g_j^{-1} \circ s_\alpha \circ g_i$
(where $s_\alpha(\mathcal{P}_i) = \mathcal{P}_j$) close
under composition to yield all $46{,}080$ stabilizer elements.

\emph{Step 2: Identify the image.}
The homomorphism
$\varphi \colon \Stab(\mathcal{P}_0) \to S_{10}$
has image of order $1{,}920$.
Every image element preserves the $5$-pair structure
(mapping perpendicular pairs to perpendicular pairs),
so $\operatorname{im}\varphi$ embeds in the group of
signed permutations on $5$ objects: each image element
permutes the $5$ pairs and optionally swaps the two
clusters within each pair.
This signed-permutation group is $\{\pm 1\}^5 \rtimes S_5
\cong W(B_5)$, of order $2^5 \cdot 5! = 3{,}840$.
The subgroup $W(D_5) = (\Z/2)^4 \rtimes S_5$ consists
of signed permutations with an \emph{even} number of
sign changes, and has order $2^4 \cdot 5! = 1{,}920$
\cite{Humphreys1990}.

Machine verification shows that every element of
$\operatorname{im}\varphi$ has an even number of
within-pair swaps.
(Explicitly: for each of the $46{,}080$ stabilizer
elements, we compute the induced signed permutation on
the $5$ pairs and verify that the product of the $5$
signs is $+1$.)
Hence $\operatorname{im}\varphi \subseteq W(D_5)$,
and the order count $1{,}920 = |W(D_5)|$ gives equality.

\emph{Step 3: Identify the kernel.}
The $24$ kernel elements (fixing every cluster setwise)
have order distribution
$\{1{:}1,\; 2{:}1,\; 3{:}8,\; 4{:}6,\; 6{:}8\}$.
Computing:
$|Z(K)| = 2$,
$[K,K] \cong Q_8$ (order $8$, quaternion group),
$K/[K,K] \cong \Z/3$.
These invariants uniquely identify $2T \cong \SL(2,\F_3)$
\cite{ConwaySloanePlesken1988}.

\emph{Step 4: Non-splitting.}
Only $2$ of $24$ kernel elements commute with all of
$\Stab(\mathcal{P}_0)$ (the identity and $-I$),
so $K$ is not central.
Lift attempts for generators of $W(D_5)$ all produce
$|K \cap H| > 1$, so no complement exists.
(A cohomological determination of the extension class in
$H^2(W(D_5), 2T)$ would sharpen this computational result;
see Open Question~1 in Section~\ref{sec:open}.)
\end{proof}

\subsection{The binary tetrahedral kernel as kinematic spinor}

\begin{remark}\label{rem:2Tspinor}
The kernel $2T \cong \SL(2,\F_3)$ is the group of $24$ unit
Hurwitz quaternions under multiplication.
Since $\HH^1 \cong \Spin(3) \cong \SU(2)$,
and $2T \subset \HH^1$,
the kernel is a \emph{finite subgroup of $\Spin(3)$}.

Each element of $2T$ acts fixed-point-freely on every $D_4$
cluster:
order-$3$ elements produce eight disjoint $3$-cycles,
order-$4$ elements produce six disjoint $4$-cycles,
and the unique order-$2$ element ($-I$) produces twelve
$2$-cycles.
These are kinematic spinors in the sense of
Definition~\ref{def:kinematic}: the orbits of each
$2T$-element on a $D_4$ shell are discrete cyclic
configurations in the spinor module of $\Spin(3)$.

The group $2T$ has spinorial character: its order-$2$ element
is $-1 \in \Spin(3)$ (the kernel of
$\Spin(3) \to \SO(3)$), which acts as the global antipodal
map on each $D_4$ shell.
\end{remark}

\begin{remark}[McKay correspondence]\label{rem:mckay}
In the McKay correspondence \cite{McKay1980}, $2T$
corresponds to the $E_6$ Dynkin diagram.
Since $E_6$ contains $\SO(10)$ as a maximal subgroup,
the exact sequence~\eqref{eq:ses} may encode the chain
$E_8 \to E_6 \to \SO(10)$ at the finite-group level.
The extension class in $H^2(W(D_5),\, 2T)$ remains to be
computed.
\end{remark}


% ============================================================
\section{The GUT Chain}
\label{sec:gut}
% ============================================================

\subsection{From $\SO(10)$ to the Standard Model}

The image of the stabilizer in $S_{10}$ is
$W(D_5) \cong (\Z/2)^4 \rtimes S_5$,
the Weyl group of $\SO(10)$.

\begin{theorem}[GUT breaking chain]\label{thm:gut}
The following Weyl-group inclusions hold:
\begin{equation}\label{eq:gutchain}
  W(D_5) \;\supset\; W(A_4) \;\supset\;
  W(A_2) \times W(A_1),
\end{equation}
with
\[
  [W(D_5) : W(A_4)] = 16 = 2^4,
  \qquad
  [W(A_4) : W(A_2)\times W(A_1)] = 10 = \binom{5}{3}.
\]
These correspond to:
\begin{enumerate}[label=(\roman*)]
\item $W(D_5) \cong W(\SO(10))$, order $1{,}920$.
\item $W(A_4) \cong S_5 \cong W(\SU(5))$, order $120$
  --- the Georgi--Glashow GUT at the Weyl level
  \cite{GeorgiGlashow1974}.
\item $W(A_2) \times W(A_1) \cong S_3 \times S_2$, order $12$
  --- the Standard Model Weyl group.
\end{enumerate}
\end{theorem}

\begin{proof}
The inclusions $W(D_n) \supset W(A_{n-1})$ and
$W(A_{n-1}) \supset W(A_k) \times W(A_{n-k-2})$ are standard
\cite{Humphreys1990,Bourbaki1968}.
We verify computationally:

\emph{$W(A_4) \subset W(D_5)$:}
The subgroup of pair permutations without within-pair swaps
is $S_5 = W(A_4)$, index
$1920/120 = 16 = |(\Z/2)^4|$.

\emph{$W(A_2)\times W(A_1) \subset W(A_4)$:}
Choosing $3$ of $5$ pairs for color and $2$ for weak gives
$S_3 \times S_2 \subset S_5$, index
$120/12 = 10 = \binom{5}{3}$.

All $10$ such choices yield conjugate copies of
$S_3 \times S_2$ inside $S_5$.
\end{proof}

\begin{remark}\label{rem:U1invisible}
The Standard Model gauge group is
$\SU(3) \times \SU(2) \times \U(1)_Y$.
The Weyl group of $\U(1)$ is trivial, so
$W(\SU(3) \times \SU(2) \times \U(1)_Y)
= S_3 \times S_2 \times \{1\} = S_3 \times S_2$.
Thus our derivation recovers the Weyl group of the full
Standard Model gauge group, but the $\U(1)_Y$ factor
is invisible at the Weyl-group level ---
it manifests only when passing from Weyl groups to
Lie algebras (where $\U(1)_Y$ appears as the trace
direction in $\mathfrak{u}(5)/\mathfrak{su}(5)$).
\end{remark}

\subsection{Ten Standard Model embeddings}

\begin{corollary}\label{cor:10embeddings}
There are exactly $10$ embeddings of the Standard Model Weyl
group $S_3 \times S_2$ inside $W(A_4) = S_5$,
all conjugate and algebraically equivalent.
\end{corollary}

The selection of a specific SM embedding requires additional
structure --- see Section~\ref{sec:discussion}.


% ============================================================
\section{Triality and Conformal Projection}
\label{sec:triality}
% ============================================================

\subsection{Two notions of triality}

The term ``triality'' is used in two distinct senses:

\begin{enumerate}[label=(\arabic*)]
\item \textbf{Mathematical triality} (the $D_4$ outer
  automorphism):
  The $D_4$ Dynkin diagram has a unique $S_3$ symmetry
  (three legs meeting at a central node),
  inducing an $S_3$ outer automorphism of $\Spin(8)$
  that permutes the vector ($8_v$), spinor ($8_s$), and
  conjugate-spinor ($8_c$) representations
  \cite{Cartan1925,Chevalley1954}.

\item \textbf{Generational triality} (the three-generation
  problem):
  The Standard Model has three fermion generations with
  identical gauge quantum numbers.
  Various authors have conjectured that mathematical triality
  explains this repetition \cite{Furey2016,Boyle2020}.
\end{enumerate}

This section provides a concrete geometric mechanism connecting
these two notions via the $4 = 1 + 3$ projection degeneracy.

\subsection{Discrete Hopf structure of the 24-cell}
\label{sec:hopfD4}

\begin{proposition}[Discrete Hopf fibers of $D_4$]
\label{prop:D4fibers}
The complex Hopf map
$h_{\C}\colon S^3 \to S^2$,
$h_{\C}(z_1,z_2) = (|z_1|^2 - |z_2|^2,\, 2z_1\bar{z}_2)$,
partitions the $24$ roots of $D_4$ into $4$ clusters of $6$
roots each.
Each cluster is a root system of type $A_2$.
This is the compositional partition $D_4 = 4 \times A_2$.
\end{proposition}

\begin{proof}
Machine-verified.
All $24$ Hopf images cluster into $4$ distinct points on $S^2$,
each with $6$ preimages spanning a $2$-dimensional subspace
and satisfying $A_2$ inner-product conditions.
\end{proof}

\begin{proposition}[Tetrahedral arrangement]
\label{prop:tetrahedron}
The $4$ Hopf images on $S^2$ form the vertices of a regular
tetrahedron: all $\binom{4}{2} = 6$ pairwise inner products
equal $-1/3$.
\end{proposition}

\begin{proof}
Machine-verified.
The value $-1/3$ is the inner product between vertices of a
regular tetrahedron inscribed in $S^2$.
\end{proof}

\subsection{The $4 = 1 + 3$ projection}

\begin{lemma}\label{lem:S3fibers}
The outer automorphism group $\Out(D_4) \cong S_3$, viewed as a
subgroup of $\mathrm{O}(4)$ via the matrices $\tau, \sigma$ of
Appendix~\ref{app:computation}, permutes the four $A_2$ Hopf
fibers of Proposition~\ref{prop:D4fibers}.
Specifically, $\tau$ (order $3$) cycles three fibers while
fixing one, and $\sigma$ (order $2$) swaps two fibers while
fixing the other two.
Conversely, any automorphism of $D_4$ that permutes the four
Hopf fibers and lies outside the Weyl group $W(D_4)$ belongs
to this $S_3$.
\end{lemma}

\begin{proof}
Machine-verified.
For each of the $4$ fibers $F_j$, we compute
$\tau(F_j)$ and $\sigma(F_j)$ and verify the resulting
permutation.
Since $W(D_4)$ acts transitively on each fiber (preserving
the partition into $4$ fibers setwise), the quotient of
the fiber-permutation group by $W(D_4) \cap \Stab$ yields
exactly $S_3 = \Out(D_4)$.
\end{proof}

\begin{proof}[Proof of Theorem~\ref{thm:D}]
By Propositions~\ref{prop:D4fibers} and~\ref{prop:tetrahedron},
the $4$ hexagonal $A_2$ fibers lie in $4$ planes whose normals
point toward the $4$ vertices of a regular tetrahedron on $S^2$.

The symmetry group of a regular tetrahedron is $S_4$,
under which all $4$ fibers are equivalent.
Projecting from $\R^4$ to $\R^3$ along one fiber normal
distinguishes that fiber: the symmetry breaks as
$S_4 \to S_3$ (stabilizer of the chosen vertex $=$ symmetry
group of the opposite triangular face).

By Lemma~\ref{lem:S3fibers}, the outer automorphism group
$\Out(D_4) \cong S_3$ acts on the four fibers by permuting
exactly the three non-distinguished fibers while fixing the
distinguished one.
Moreover, $S_3$ fixes the $2$-dimensional plane $V_0$ of the
distinguished fiber pointwise (since $V_0$ is the fiber that
$S_3$ stabilizes), and therefore also fixes the orthogonal
complement $V_0^\perp$ pointwise.
Consequently, the projection $\pi_{\mathbf{n}}$ commutes
with $S_3$ for \emph{every} choice of
$\mathbf{n} \in V_0^\perp$, and the residual $S_3$
symmetry of the projected configuration is independent
of this choice.
\end{proof}

\subsection{Semi-conformal angle preservation}
\label{sec:conformal}

We now prove that the orthogonal projection from $\R^4$ to
$\R^3$ preserves the root-system angular structure on the
outer shell --- the key geometric claim that gives the
$D_4$ angle set a faithful $3$-dimensional representative.

\begin{proof}[Proof of Theorem~\ref{thm:E}]
Machine-verified.
Let $P = \{F_0, F_1, F_2, F_3\}$ be a compositional partition
of $D_4$ into four $A_2$ fibers
(Proposition~\ref{prop:D4fibers}).
The fiber $F_0$ spans a $2$-dimensional subspace
$V_0 \subset \R^4$, whose orthogonal complement $V_0^\perp$
is also $2$-dimensional.
Let $\mathbf{n}$ be any unit vector in $V_0^\perp$.
The orthogonal projection
$\pi_{\mathbf{n}} \colon \R^4 \to \mathbf{n}^\perp \cong \R^3$
sends each root $r \mapsto r - \langle r, \mathbf{n}\rangle
\,\mathbf{n}$.
(The result below is verified for all directions $\mathbf{n}$
in $V_0^\perp$; the shell decomposition and angle sets are
independent of the choice of $\mathbf{n}$ within $V_0^\perp$.)

\emph{Step 1: Shell decomposition.}
Under $\pi_{\mathbf{n}}$, the $24$ roots separate into
two concentric spherical shells by $3$-dimensional radius:
\begin{itemize}
\item \emph{Outer shell} (radius $\sqrt{2}$):
  $12$ roots project to $12$ distinct points.
\item \emph{Inner shell} (radius $1$):
  $12$ roots project to $6$ distinct positions
  (pairs of roots from the same fiber collapse under projection).
\end{itemize}

\emph{Step 2: Outer shell = cuboctahedron.}
The $12$ outer-shell points form a \emph{cuboctahedron}:
the root system $A_3$ (equivalently $D_3$), the unique
rank-$3$ root system with $12$ roots.
We verify this by checking:
(i) all $12$ projected vectors have equal norm $\sqrt{2}$;
(ii) the Gram matrix of pairwise inner products matches
the $A_3$ root system; and
(iii) all $\binom{12}{2} = 66$ pairwise angles lie in
$\{60^\circ, 90^\circ, 120^\circ, 180^\circ\}$.

This angle set $\{60^\circ, 90^\circ, 120^\circ, 180^\circ\}$
is exactly the angle set of $D_4$ itself: the inner products
between $D_4$ roots of squared norm $2$ are
$\langle \alpha, \beta \rangle \in \{-2, -1, 0, 1\}$,
giving $\cos\theta \in \{-1, -1/2, 0, 1/2\}$,
i.e.\ $\theta \in \{180^\circ, 120^\circ, 90^\circ, 60^\circ\}$.

\emph{Step 3: Inner shell = octahedron.}
The $6$ distinct inner-shell positions form a regular
octahedron with pairwise angles in $\{90^\circ, 180^\circ\}$
--- the $D_4$ angle subset corresponding to perpendicular
and anti-parallel roots.

\emph{Step 4: Cross-shell angles.}
Between the outer and inner shells, additional angles
$45^\circ$ and $135^\circ$ appear.
These are \emph{not} root-system angles.
\end{proof}

\begin{remark}[Semi-conformal interpretation]
\label{rem:semiconformal}
The cuboctahedral outer shell is a faithful
$3$-dimensional representative of $D_4$'s angular
structure: it preserves the full angle set
$\{60^\circ, 90^\circ, 120^\circ, 180^\circ\}$ without
introducing extraneous angles.
The inner octahedral shell preserves a subset of these angles.
The qualifier ``semi-conformal'' reflects the fact that
while each shell individually preserves $D_4$-compatible
angles, the cross-shell angles introduce values ($45^\circ$
and $135^\circ$) that lie outside the root-system angle set.

The cuboctahedral shell provides a concrete
$3$-dimensional realization of $D_4$'s angular structure
that may be useful for visualization and for connecting
to $3$-dimensional physical models.
\end{remark}

\begin{remark}[Kinematic spinor perspective on projection]
\label{rem:projectionspinor}
The projection from $\R^4$ to $\R^3$ is, from the spinor
perspective, the passage from the Level~2 kinematic spinor
($\Spin(4)$ acting on $D_4$) to its $\Spin(3)$ shadow.
The distinguished fiber $F_0$ is the ``gauge'' choice:
it determines which $\U(1) \subset \Spin(4)$ is projected out,
breaking $\Spin(4) \cong \SU(2)_L \times \SU(2)_R$ to the
diagonal $\Spin(3) \cong \SU(2)$.
The three surviving fibers carry the $S_3$ triality,
and the cuboctahedral shell is the geometric form of
$D_4$'s angular information as seen by $\Spin(3)$ observers.
\end{remark}

\subsection{The Pauli matrix analogy}

The $4 = 1 + 3$ split has a natural algebraic interpretation.
The $2 \times 2$ Hermitian matrices decompose as
$H = t\, I + x\, \sigma_1 + y\, \sigma_2 + z\, \sigma_3$,
where $\{I, \sigma_1, \sigma_2, \sigma_3\}$ are the Pauli
matrices.
The identity $I$ is the trace (``$1$'') and
$\{\sigma_1, \sigma_2, \sigma_3\}$ are the traceless
directions (``$3$''), transforming as a vector under $\SO(3)$.

In our framework:
\begin{itemize}
\item The chosen projection fiber corresponds to $I$.
\item The remaining $3$ fibers correspond to
  $\{\sigma_1, \sigma_2, \sigma_3\}$ (triplet under $S_3$).
\item The choice of fiber is a symmetry-breaking:
  all $4$ choices are equivalent \emph{a priori}.
\end{itemize}
This decomposition is structurally identical to the passage
from quaternionic structure ($\dim_\R = 4$) to complexified
vector structure ($1 + 3$) that underlies $\SU(2)$
representation theory.


% ============================================================
\section{Discussion}
\label{sec:discussion}
% ============================================================

\subsection{Summary of results}

Starting from the four normed division algebras and their
integral lattices, using only discrete cyclic group actions
(kinematic spinors) that preserve these lattices, we have
derived:
\begin{enumerate}
\item The \emph{compositional hierarchy}
  $A_1 \subset A_2 \subset D_4 \subset E_8$,
  where each root system is a compound of oriented copies
  of the previous one, generated by kinematic spinors in
  the corresponding spin group (Theorem~\ref{thm:A}).
\item The \emph{GUT group $\SO(10)$} as the Weyl-group-level
  structure of the $E_8$ partition stabilizer, with the
  binary tetrahedral group $2T$ as a kinematic-spinor kernel
  (Theorem~\ref{thm:B}).
\item The \emph{Standard Model Weyl group}
  $W(A_2) \times W(A_1) \cong S_3 \times S_2$ from the GUT breaking chain,
  with the correct Weyl group at each stage (Theorem~\ref{thm:C}).
  By the reconstruction theorem for simply-laced root systems
  (Remark~\ref{rem:reconstruction}), this canonically determines
  the continuous semisimple factor $\SU(3)\times\SU(2)$;
  only $\U(1)_Y$ and the $\Z_6$ quotient remain underived.
\item The \emph{three-fold degeneracy} $4 = 1 + 3$ from the
  projection of $D_4$ Hopf fibers, with residual $S_3 = \Out(D_4)$
  (Theorem~\ref{thm:D}).
\item The \emph{semi-conformal angle preservation}: the
  cuboctahedral shell in $\R^3$ faithfully represents $D_4$'s
  full angle set (Theorem~\ref{thm:E}).
\end{enumerate}

\subsection{What is not derived}

This paper derives the \emph{Weyl group} of the Standard
Model gauge group, not the continuous gauge group itself
or its dynamics:
\begin{itemize}
\item For the semisimple sector, the Weyl group
  \emph{does} canonically reconstruct the continuous Lie
  group $\SU(3)\times\SU(2)$
  (Remark~\ref{rem:reconstruction}).
  What remains beyond reach of the Weyl group:
  the abelian factor $\U(1)_Y$ (trivial Weyl group),
  the $\Z_6$ center quotient, and root lengths in
  non-simply-laced types.
\item No Lagrangian or equations of motion.
\item No fermion masses or mixing matrices.
\item No Higgs mechanism or symmetry-breaking potential.
\item No specific representation content (though the $5+5$
  bipartite structure suggests $\mathbf{10}\oplus
  \overline{\mathbf{5}}$ of $\SU(5)$).
\end{itemize}

\begin{remark}[Reconstruction of continuous Lie groups]
\label{rem:reconstruction}
A classical result in the theory of root systems asserts that,
for simply-laced (ADE) types, the Weyl group $W(\Phi)$
uniquely determines the root system $\Phi$, and hence---via
the Serre presentation---the semisimple Lie algebra
$\mathfrak{g}$ and its compact simply-connected Lie group $G$
\cite{Humphreys1990, Bourbaki1968}.
The chain is:
\[
  W(\Phi)
  \;\xrightarrow{\;\text{Coxeter = Dynkin}\;}
  \text{diagram}
  \;\xrightarrow{\;\text{Cartan matrix}\;}
  \Phi
  \;\xrightarrow{\;\text{Serre relations}\;}
  \mathfrak{g}
  \;\xrightarrow{\;\exp\;}
  G.
\]
For non-simply-laced types this fails ($B_n$ and $C_n$ share
the same Weyl group), but the types $A_n$ appearing in our
derivation are simply-laced.

Theorem~\ref{thm:C} identifies
$W(A_2) \times W(A_1) \cong S_3 \times S_2$
through the ADE inclusion chain
$W(D_5) \supset W(A_4) \supset W(A_2) \times W(A_1)$,
so each factor carries a canonical ADE label.
Applying the reconstruction:
\[
  W(A_2) \;\longrightarrow\; A_2
  \;\longrightarrow\; \mathfrak{su}(3)
  \;\longrightarrow\; \SU(3),
  \qquad
  W(A_1) \;\longrightarrow\; A_1
  \;\longrightarrow\; \mathfrak{su}(2)
  \;\longrightarrow\; \SU(2).
\]
The passage from discrete Weyl group to continuous Lie group
is therefore \emph{canonical} for the semisimple sector.
The only Standard Model data not captured at the Weyl-group
level are the abelian factor $\U(1)_Y$ (whose Weyl group is
trivial) and the $\Z_6$ center quotient (invisible to Weyl
reflections).
\end{remark}

\subsection{The selection problem}\label{sec:selection}

The $10$ equivalent SM embeddings
(Corollary~\ref{cor:10embeddings})
correspond to choosing $3$ of $5$ perpendicular pairs for color
and $2$ for weak isospin.
The $4 = 1 + 3$ triality projection
(Theorem~\ref{thm:D}) applies \emph{locally} to each $D_4$
cluster, splitting its four $A_2$ fibers into a distinguished
singlet and a triplet.
Whether this local fiber triality can induce a \emph{global}
selection among the five perpendicular $D_4$ pairs in $E_8$
--- and thereby single out a unique SM embedding --- remains
an open problem.
A complete proof would require constructing an explicit map
from the fiber-selection data within each $D_4$ cluster to a
canonical choice of a $3$-subset of the $5$ pair-classes in
the global $W(D_5)$ structure.

\subsection{The kinematic spinor hierarchy in perspective}

The nested structure of kinematic spinors described in this
paper --- $\Spin(2)$ generating $A_2$ from $A_1$,
$\U(1) \subset \Spin(4)$ organizing $D_4$ from $A_2$,
$2T \subset \Spin(3)$ acting on $D_4$ fibers in $E_8$ ---
mirrors the nested Hopf fibrations:
\[
  S^1 \to S^0 \quad(\R),\qquad
  S^3 \to S^2 \quad(\C),\qquad
  S^7 \to S^4 \quad(\HH).
\]
Hurwitz's theorem \cite{Hurwitz1898} ensures no
$16$-dimensional normed division algebra exists, so there is
no further division-algebra root lattice beyond $E_8$ on which
the hierarchy could continue.
(The octonionic Hopf fibration $S^{15} \to S^8$ exists
but has no higher-dimensional division-algebra lattice to act upon.)
The Standard Model Weyl group emerges at the top level not
because we have chosen it but because the rigid geometry of
$E_8$, constrained by the division-algebra lattice structure
and the finite Hopf-fiber symmetries, naturally yields a
canonical $S_3 \times S_2$ structure through the GUT chain.

From this perspective, the kinematic spinor is the discrete
geometric mechanism that translates the topological constraint
(existence of Hopf fibrations in and only in dimensions
$1, 2, 4, 8$) into the algebraic structure (the Weyl group
of the Standard Model gauge group).

\subsection{Open questions}\label{sec:open}

\begin{enumerate}
\item \textbf{Extension classification.}
  The non-split extension~\eqref{eq:ses}, with $2T$ corresponding
  to $E_6$ via McKay, may encode $E_8 \to E_6 \to \SO(10)$
  at the finite-group level.
  The extension class in $H^2(W(D_5), 2T)$ remains to be
  computed.

\item \textbf{Representation content.}
  The $5+5$ bipartite structure of the partition
  (Proposition~\ref{prop:perppairs}) is suggestive of the
  $\mathbf{10} \oplus \overline{\mathbf{5}}$ decomposition
  of $\SU(5)$ representations.
  Extending from Weyl-group level to full Lie-algebra
  representations --- identifying the $\mathbf{16}$-plet of
  $\SO(10)$ with lattice structures --- is a natural next step.

\item \textbf{Physical subscripts.}
  While we identify $\SU(3)$ with the triality triplet
  and $\SU(2)$ with the weak doublet, a derivation of
  the physical quantum numbers (color, chirality, hypercharge)
  from lattice geometry is not completed here.

\item \textbf{Conformal projection to lower dimensions.}
  The semi-conformal result (Theorem~\ref{thm:E}) shows that
  shell-by-shell angle preservation holds for $D_4 \to \R^3$.
  Whether an analogous statement holds for $E_8 \to \R^3$
  (through the full $\Spin(8) \to \Spin(3)$ reduction) is open.
\end{enumerate}

\subsection{Relation to prior work}

Section~\ref{sec:priorapproaches} surveys the main research
programs connecting division algebras and exceptional structures
to the Standard Model.
Here we note specific points of contact and novelty.

Dechant's spinor induction \cite{Dechant2016} is the closest
precursor: both approaches use Clifford-algebraic actions on
root-system lattices.
Our construction differs in using a \emph{recursive},
multi-level hierarchy ($A_1 \to A_2 \to D_4 \to E_8$) and
in extracting the Standard Model Weyl group --- and hence,
by the reconstruction theorem
(Remark~\ref{rem:reconstruction}), the continuous semisimple
factor $\SU(3)\times\SU(2)$ --- from the Hopf partition
stabilizer at the top level.

Theorem~\ref{thm:A}(c) (the $E_8$ partition) builds on the
Hopf construction of \cite{ClawsonFangIrwin2024}.
The compositional hierarchy as a three-level spinor
structure (Theorem~\ref{thm:A}),
the identification of the Hopf partition stabilizer as a
non-split $2T$-extension of $W(D_5)$ (Theorem~\ref{thm:B}),
the GUT chain extraction (Theorem~\ref{thm:C}),
and the semi-conformal angle preservation
(Theorem~\ref{thm:E}) appear to be new.


% ============================================================
% APPENDIX
% ============================================================
\appendix

\section{Computational Methods}\label{app:computation}

All machine verifications use exact integer and rational
arithmetic on the $E_8$ root system in standard coordinates:
$112$ integer-type roots $\pm e_i \pm e_j$ ($i < j$) plus
$128$ half-integer-type roots
$(\pm\tfrac{1}{2},\ldots,\pm\tfrac{1}{2})$ with an even
number of minus signs.
Group elements are stored as permutations of the $240$ roots;
all operations are exact.

The triality matrices for $D_4$ are:
\[
  \tau = \frac{1}{2}
  \begin{pmatrix}
    1 & 1 & 1 & 1 \\
    1 & 1 & -1 & -1 \\
    1 & -1 & 1 & -1 \\
    -1 & 1 & 1 & -1
  \end{pmatrix},
  \qquad
  \sigma\colon (x_1, x_2, x_3, x_4) \mapsto (x_1, x_2, x_4, x_3).
\]
Here $\tau$ has order $3$ (cycling the three legs of the $D_4$
Dynkin diagram) and $\sigma$ has order $2$ (swapping two legs).
Note that $\tau$ is \emph{not} a signed permutation matrix:
it is an outer automorphism of $D_4$, not a Weyl group element.

\subsection{Machine-verified results}

\begin{center}
\begin{tabular}{@{}llr@{}}
\toprule
Result & Method & Value \\
\midrule
$|E_8|$ & Enumeration & $240 = 112 + 128$ \\
Hopf clusters & Quaternionic Hopf map & 10 of 24 \\
$D_4$ verification & Cartan matrix & All 10 confirmed \\
Perpendicular pairs & Cross-Gram matrix & 5 pairs \\
Partition orbit (BFS) & 120 generators & 15{,}120 \\
$|\Stab|$ & Orbit-stabilizer & 46{,}080 \\
$|\ker\varphi|$ & Permutation kernel & 24 \\
$|\operatorname{im}\varphi|$ & Image in $S_{10}$ & 1{,}920 \\
$\ker\varphi \cong 2T$ & Invariants & $\SL(2,\F_3)$ \\
$\operatorname{im}\varphi \cong W(D_5)$ & Pair preservation
  & $(\Z/2)^4 \rtimes S_5$ \\
Non-split & Centrality + complement & Confirmed \\
GUT chain indices & Subgroup construction & 16, 10 \\
$D_4 = 4 \times A_2$ & Exhaustive & 8 partitions \\
Triality orbits & $S_3$ action & $2 + 6$ \\
Tetrahedral arrangement & Inner products & All $= -1/3$ \\
Cuboctahedral shell & Projection + angles & $\{60,90,120,180\}^\circ$ \\
Octahedral shell & Projection + angles & $\{90, 180\}^\circ$ \\
\bottomrule
\end{tabular}
\end{center}

Source code for all computations is publicly available at\\
\url{https://github.com/QGRKlee/CCT-StandardModel}\\
and will be archived with a DOI prior to publication.
Below we provide algorithmic specifications sufficient
for independent reimplementation.

\subsection{Data representations}

\emph{Roots.}
Each $E_8$ root is stored as an $8$-tuple of integers or
half-integers.
The $240$ roots are indexed $\{1,\ldots,240\}$; each root
$r_i$ also stores $-r_i$'s index for $O(1)$ negation.

\emph{Partitions.}
A partition $\mathcal{P}$ is represented as a list of $10$
clusters, each cluster a sorted tuple of root indices.
The canonical form is obtained by sorting the $10$ clusters
lexicographically.
Two partitions are equal if and only if their canonical
forms are identical (no hash collisions possible).

\emph{Reflections.}
The Weyl reflection $s_\alpha$ through root $\alpha$
(Equation~\ref{eq:reflection}) acts on a partition by
applying $s_\alpha$ to each root in each cluster, re-sorting
each cluster, and re-canonicalizing.
Since $s_\alpha(r) = r - \langle r,\alpha\rangle\,\alpha$
with $\langle\alpha,\alpha\rangle = 2$, all arithmetic is
exact (integer or half-integer).

\subsection{BFS orbit enumeration (Theorem~\ref{thm:orbit})}

\begin{enumerate}
\item Initialize queue $Q \gets \{\mathcal{P}_0\}$,
  visited set $V \gets \{\text{canonical}(\mathcal{P}_0)\}$.
\item For each partition $\mathcal{P}$ dequeued from $Q$:
  for each of the $120$ root reflections
  $s_\alpha$ (one per antipodal pair):
  \begin{enumerate}
  \item Compute $\mathcal{P}' = s_\alpha(\mathcal{P})$
    and its canonical form $c'$.
  \item If $c' \notin V$: add $c'$ to $V$, enqueue
    $\mathcal{P}'$, and record the coset representative
    $g_{\mathcal{P}'} = s_\alpha \circ g_{\mathcal{P}}$.
  \end{enumerate}
\item Terminate when $Q$ is empty.
  Since $W(E_8)$ is generated by root reflections and the
  orbit is finite, BFS is guaranteed to visit every element
  of the orbit.
  Result: $|V| = 15{,}120$.
\end{enumerate}

\subsection{Stabilizer construction (Theorem~\ref{thm:stab})}

Schreier generators \cite{Sims1970}: for each reflection
$s_\alpha$ and each visited partition $\mathcal{P}_i$ with
$s_\alpha(\mathcal{P}_i) = \mathcal{P}_j$, the element
$g_j^{-1} \circ s_\alpha \circ g_i$ stabilizes
$\mathcal{P}_0$.
Collecting and closing under composition yields all
$46{,}080$ stabilizer elements stored as permutations
of the $240$ root indices.

\subsection{Kernel identification}

The homomorphism $\varphi$ records each stabilizer element's
induced permutation on the $10$ cluster labels.
The kernel ($24$ elements with trivial cluster permutation)
is identified by computing:
order distribution
$\{1{:}1,\; 2{:}1,\; 3{:}8,\; 4{:}6,\; 6{:}8\}$,
center order $|Z(K)| = 2$,
and derived subgroup $[K,K] \cong Q_8$ ---
invariants that uniquely characterize
$2T \cong \SL(2,\F_3)$ \cite{ConwaySloanePlesken1988}.


% ============================================================
% REFERENCES
% ============================================================
\begin{thebibliography}{99}

\bibitem{Adams1960}
J.\,F.~Adams,
\emph{On the non-existence of elements of Hopf invariant one},
Ann.\ of Math.\ \textbf{72} (1960), 20--104.

\bibitem{Baez2002}
J.\,C.~Baez,
\emph{The octonions},
Bull.\ Amer.\ Math.\ Soc.\ \textbf{39} (2002), 145--205.

\bibitem{Bourbaki1968}
N.~Bourbaki,
\emph{Groupes et alg\`ebres de Lie},
Chap.\ 4--6, Hermann, Paris, 1968.

\bibitem{Boyle2020}
L.~Boyle,
\emph{The Standard Model, the exceptional Jordan algebra,
and triality},
arXiv:2006.16265, 2020.

\bibitem{Cartan1925}
E.~Cartan,
\emph{Le principe de dualit\'e et la th\'eorie des groupes
simples et semi-simples},
Bull.\ Sci.\ Math.\ \textbf{49} (1925), 361--374.

\bibitem{Chevalley1954}
C.~Chevalley,
\emph{The algebraic theory of spinors},
Columbia Univ.\ Press, 1954.

\bibitem{ClawsonFangIrwin2024}
R.~Clawson, F.~Fang, and K.~Irwin,
\emph{From the Fibonacci Icosagrid to $E_8$ (Part~II)},
Symmetry, 2024.

\bibitem{ConwaySloane1999}
J.\,H.~Conway and N.\,J.\,A.~Sloane,
\emph{Sphere Packings, Lattices and Groups},
3rd ed., Springer, 1999.

\bibitem{ConwaySloanePlesken1988}
J.\,H.~Conway, R.\,T.~Curtis, S.\,P.~Norton,
R.\,A.~Parker, and R.\,A.~Wilson,
\emph{ATLAS of Finite Groups},
Oxford Univ.\ Press, 1985.

\bibitem{Coxeter1946}
H.\,S.\,M.~Coxeter,
\emph{Integral Cayley numbers},
Duke Math.\ J.\ \textbf{13} (1946), 561--578.

\bibitem{Dechant2016}
P.-P.~Dechant,
\emph{The birth of $E_8$ out of the spinors of the icosahedron},
Proc.\ R.\ Soc.\ A \textbf{472} (2016), 20150504.

\bibitem{DistlerGaribaldi2009}
J.~Distler and S.~Garibaldi,
\emph{There is no ``Theory of Everything'' inside $E_8$},
Comm.\ Math.\ Phys.\ \textbf{298} (2010), 419--436.

\bibitem{Dixon1994}
G.\,M.~Dixon,
\emph{Division Algebras: Octonions, Quaternions, Complex
Numbers and the Algebraic Design of Physics},
Kluwer, 1994.

\bibitem{FangIrwin2024}
F.~Fang and K.~Irwin,
\emph{From the Fibonacci Icosagrid to $E_8$ (Part~I)},
Symmetry, 2024.

\bibitem{FritzschMinkowski1975}
H.~Fritzsch and P.~Minkowski,
\emph{Unified interactions of leptons and hadrons},
Ann.\ Phys.\ \textbf{93} (1975), 193--266.

\bibitem{Furey2016}
C.~Furey,
\emph{Standard Model physics from an algebra?},
Ph.D.\ thesis, University of Waterloo, 2016.

\bibitem{GeorgiGlashow1974}
H.~Georgi and S.\,L.~Glashow,
\emph{Unity of all elementary-particle forces},
Phys.\ Rev.\ Lett.\ \textbf{32} (1974), 438--441.

\bibitem{Hopf1935}
H.~Hopf,
\emph{\"Uber die Abbildungen von Sph\"aren auf Sph\"aren
niedrigerer Dimension},
Fund.\ Math.\ \textbf{25} (1935), 427--440.

\bibitem{Hurwitz1898}
A.~Hurwitz,
\emph{\"Uber die Composition der quadratischen Formen von
beliebig vielen Variablen},
Nachr.\ Ges.\ Wiss.\ G\"ottingen (1898), 309--316.

\bibitem{Humphreys1990}
J.\,E.~Humphreys,
\emph{Reflection Groups and Coxeter Groups},
Cambridge Univ.\ Press, 1990.

\bibitem{Krasnov2019}
K.~Krasnov,
\emph{$\SO(9)$ characterisation of the Standard Model gauge group},
J.\ Math.\ Phys.\ \textbf{62} (2021), 021703;
arXiv:1912.11282, 2019.

\bibitem{LawsonMichelsohn1989}
H.\,B.~Lawson and M.-L.~Michelsohn,
\emph{Spin Geometry},
Princeton Univ.\ Press, 1989.

\bibitem{Lisi2007}
A.\,G.~Lisi,
\emph{An Exceptionally Simple Theory of Everything},
arXiv:0711.0770, 2007.

\bibitem{McKay1980}
J.~McKay,
\emph{Graphs, singularities, and finite groups},
Proc.\ Sympos.\ Pure Math.\ \textbf{37} (1980), 183--186.

\bibitem{Sims1970}
C.\,C.~Sims,
\emph{Computational methods in the study of permutation groups},
in: Computational Problems in Abstract Algebra,
Pergamon, 1970, pp.\ 169--183.

\bibitem{TodorovDuboisViolette2018}
I.\,T.~Todorov and M.~Dubois-Violette,
\emph{Deducing the symmetry of the Standard Model from the
automorphism and structure groups of the exceptional Jordan
algebra},
Int.\ J.\ Mod.\ Phys.\ A \textbf{33} (2018), 1850118.

\bibitem{Viazovska2017}
M.\,S.~Viazovska,
\emph{The sphere packing problem in dimension $8$},
Ann.\ of Math.\ \textbf{185} (2017), 991--1015.

\end{thebibliography}

\end{document}
